\chapter{Cilium e eBPF}
\label{cap:cilium}

O Cilium é um projeto de código aberto que implementa rede, segurança e
observabilidade para Kubernetes usando eBPF como base \cite{ciliumdocs}. Ele é o
exemplo mais completo do que o capítulo \ref{cap:k8s} descreveu em termos gerais: um
plugin CNI que anexa programas eBPF a cada \textit{pod}, substitui o
\texttt{kube-proxy}, aplica políticas por identidade e expõe os fluxos observados por
meio do Hubble \cite{hubbledocs}.

\textbf{Aviso de escopo.} Assim como no capítulo anterior, este capítulo é
\textbf{conceitual}. Não havia cluster Kubernetes, \texttt{kubectl}, \texttt{kind} nem
a CLI do Cilium na máquina em que o manual foi escrito. Nenhum comando deste capítulo
foi executado; as formas e as opções foram tiradas da documentação oficial
\cite{ciliumdocs, hubbledocs} e podem mudar entre versões. Antes de usar, confira a
documentação da versão instalada. Os manifestos YAML foram verificados apenas quanto à
sintaxe.

% ----------------------------------------------------------------
\section{Cilium como CNI}
% ----------------------------------------------------------------

\subsection*{Componentes}

A arquitetura do Cilium tem poucas peças \cite{ciliumdocs}:

\begin{itemize}
  \item o \textbf{agente} (\texttt{cilium-agent}), executado como \texttt{DaemonSet},
        isto é, um por nó. Ele recebe do Kubernetes os \textit{pods}, serviços e
        políticas, compila e carrega os programas eBPF e mantém os mapas BPF;
  \item o \textbf{plugin CNI}, invocado pelo \textit{kubelet} a cada criação de
        \textit{pod}, que aciona o agente (conforme a operação \texttt{ADD} da seção
        CNI do capítulo \ref{cap:k8s});
  \item o \textbf{operador}, responsável por tarefas de escopo de cluster, como a
        alocação de blocos de endereços;
  \item o \textbf{Hubble}, camada de observabilidade, apresentada adiante.
\end{itemize}

\subsection*{Instalação}

A forma mais direta de instalar o Cilium em um cluster novo é a CLI do projeto. O
bloco a seguir é um \textbf{comando de cluster, não executado}:

\begin{lstlisting}[language=bash, caption={Instalação e verificação do Cilium com a CLI (comando de cluster, não executado).}, label=lst:cilium-install]
cilium install
cilium status --wait
\end{lstlisting}

Alternativamente, o Cilium é instalado com Helm, o que permite fixar a versão e ativar
funcionalidades por parâmetros. Os valores abaixo ativam a substituição do
\texttt{kube-proxy} e o Hubble; os nomes de parâmetros mudam entre versões, e o
\texttt{kubeProxyReplacement} exige que o cluster tenha sido criado sem o
\texttt{kube-proxy} ou que ele seja removido:

\begin{lstlisting}[language=bash, caption={Instalação via Helm com Hubble e substituição do kube-proxy (comando de cluster, não executado).}, label=lst:cilium-helm]
helm repo add cilium https://helm.cilium.io/
helm install cilium cilium/cilium --namespace kube-system \
     --set kubeProxyReplacement=true \
     --set hubble.relay.enabled=true \
     --set hubble.ui.enabled=true
\end{lstlisting}

\subsection*{O que o Cilium carrega no kernel}

Em cada nó, o agente anexa programas eBPF principalmente nos pontos estudados nos
capítulos \ref{cap:xdp} e \ref{cap:tc}: TC (\textit{ingress} e \textit{egress}) nas
interfaces dos \textit{pods} e do nó, XDP opcionalmente na interface externa para
balanceamento, e programas de \textit{cgroup} para a tradução de serviços no
\texttt{connect()} (Listing \ref{lst:k8s-connect4}) \cite{ciliumdocs}.

Esses programas são visíveis com as mesmas ferramentas do restante do manual. Os
comandos a seguir são \textbf{executáveis em um nó} (com privilégio de administrador),
mas só produzem resultado onde o Cilium estiver instalado; a saída, portanto, não é
reproduzida:

\begin{lstlisting}[language=bash, caption={Inspeção dos programas do Cilium em um nó com bpftool (requer root em um nó com Cilium).}, label=lst:cilium-bpftool]
sudo bpftool prog show | grep -i cil
sudo bpftool net show
sudo bpftool map show | head
\end{lstlisting}

O \texttt{bpftool net show} lista os programas XDP e TC anexados a cada interface. Em
um nó com Cilium, espera-se encontrar programas TC nas interfaces dos \textit{pods}
(nomes iniciados por \texttt{lxc}) e na interface do nó; a ferramenta confirma na
prática o que a seção anterior descreveu, e é um bom exercício comparar o que ela
lista com o que o capítulo \ref{cap:tc} ensinou a anexar manualmente.

% ----------------------------------------------------------------
\section{Network Policies}
% ----------------------------------------------------------------

O Cilium aceita a \texttt{NetworkPolicy} padrão do Kubernetes (Listing
\ref{lst:k8s-netpol}) e oferece o recurso próprio \texttt{CiliumNetworkPolicy}, mais
expressivo \cite{ciliumdocs}. A diferença central está na forma como as políticas são
aplicadas. O Cilium atribui a cada \textit{pod} uma \textbf{identidade de segurança},
derivada do conjunto de rótulos relevantes, e a identidade viaja com o pacote
(por exemplo, em um campo do encapsulamento ou consultada em mapa no destino). A regra
de política, no fim, é uma entrada de mapa BPF do tipo \textit{(identidade de origem,
porta, protocolo) $\rightarrow$ permitir}, e não uma lista de endereços IP.

Duas consequências práticas merecem destaque: a política acompanha o \textit{pod} quando
ele é recriado com outro endereço, e a verificação tem custo praticamente
independente do número de \textit{pods} \cite{ciliumdocs}.

O manifesto a seguir é um \textbf{exemplo de manifesto}, \textbf{não aplicado}. Ele
equivale ao Listing \ref{lst:k8s-netpol}, mas acrescenta uma regra de camada 7:
somente requisições \texttt{GET} em caminhos que começam por \texttt{/api/}.

\begin{lstlisting}[caption={CiliumNetworkPolicy com regra HTTP na camada 7 (exemplo de manifesto).}, label=lst:cilium-cnp]
apiVersion: cilium.io/v2
kind: CiliumNetworkPolicy
metadata:
  name: api-somente-get
  namespace: loja
spec:
  endpointSelector:
    matchLabels:
      app: api
  ingress:
    - fromEndpoints:
        - matchLabels:
            app: web
      toPorts:
        - ports:
            - port: "8080"
              protocol: TCP
          rules:
            http:
              - method: "GET"
                path: "/api/.*"
\end{lstlisting}

A regra de camada 7 tem um custo que a de camada 3/4 não tem: o Cilium precisa
\textbf{interpretar o protocolo de aplicação}, e para isso redireciona o tráfego
afetado a um \textit{proxy} (baseado em Envoy) em vez de decidir apenas no programa
eBPF \cite{ciliumdocs}. É uma troca consciente entre expressividade e desempenho:
use a regra de camada 7 apenas onde a distinção por método ou caminho for realmente
necessária.

Também é possível controlar a \textbf{saída} (\textit{egress}) por nome de domínio,
útil para permitir que uma aplicação acesse um serviço externo específico sem
manter listas de IP. O Cilium resolve essa regra observando as respostas DNS do próprio
\textit{pod} \cite{ciliumdocs}:

\begin{lstlisting}[caption={Política de saída por nome de domínio (exemplo de manifesto).}, label=lst:cilium-fqdn]
apiVersion: cilium.io/v2
kind: CiliumNetworkPolicy
metadata:
  name: web-saida-dns
  namespace: loja
spec:
  endpointSelector:
    matchLabels:
      app: web
  egress:
    - toEndpoints:
        - matchLabels:
            k8s:io.kubernetes.pod.namespace: kube-system
            k8s:k8s-app: kube-dns
      toPorts:
        - ports:
            - port: "53"
              protocol: ANY
          rules:
            dns:
              - matchPattern: "*"
    - toFQDNs:
        - matchName: "api.exemplo.com.br"
      toPorts:
        - ports:
            - port: "443"
              protocol: TCP
\end{lstlisting}

A primeira regra de saída libera as consultas ao DNS do cluster, que é pré-requisito
para a regra \texttt{toFQDNs} funcionar, já que é a observação dessas respostas que
informa ao Cilium quais endereços correspondem ao nome. O domínio
\texttt{api.exemplo.com.br} é fictício.

Aplicar e listar os recursos usa o \texttt{kubectl}:

\begin{lstlisting}[language=bash, caption={Aplicação das políticas do Cilium (comando de cluster, não executado).}, label=lst:cilium-apply]
kubectl apply -f api-somente-get.yaml
kubectl get cnp -n loja
\end{lstlisting}

% ----------------------------------------------------------------
\section{Observabilidade com Hubble}
% ----------------------------------------------------------------

O Hubble é a camada de observabilidade do Cilium. Ele aproveita o fato de que os
programas eBPF do Cilium já veem cada pacote para emitir \textbf{eventos de fluxo},
que incluem origem e destino por identidade (nome do \textit{pod}, \textit{namespace}),
protocolo, porta e, o ponto mais valioso, o \textbf{veredito} da política
(\texttt{FORWARDED}, \texttt{DROPPED}) \cite{hubbledocs}. Isso responde à pergunta
que o capítulo \ref{cap:k8s} deixou aberta: qual política derrubou a conexão?

Cada agente mantém um histórico limitado dos fluxos do seu nó. O componente
\textit{Hubble Relay} agrega os agentes e oferece uma única API para o cluster inteiro;
a CLI \texttt{hubble} e a interface web (Hubble UI) consomem essa API
\cite{hubbledocs}. O bloco a seguir ativa e usa o Hubble. É um \textbf{comando de
cluster, não executado}:

\begin{lstlisting}[language=bash, caption={Ativação e consulta de fluxos com o Hubble (comando de cluster, não executado).}, label=lst:hubble-observe]
cilium hubble enable --ui
cilium hubble port-forward &

hubble status
hubble observe --namespace loja
hubble observe --namespace loja --verdict DROPPED --follow
hubble observe --to-pod loja/api --protocol http
cilium hubble ui
\end{lstlisting}

O segundo comando de \texttt{hubble observe} é o mais útil para depuração de
política: ele acompanha em tempo real apenas os fluxos descartados. Cada linha traz a
origem, o destino, o protocolo e o veredito, de modo que o \textit{pod} bloqueado e a
política responsável ficam evidentes. A saída é por \textit{pod} (nomes de recurso do
Kubernetes) e não por endereço IP, o que a torna legível em um ambiente onde os
endereços mudam o tempo todo. Neste manual, não há saída real a mostrar.

O Hubble também gera métricas no formato Prometheus (contagem de fluxos por
veredito, consultas DNS e requisições HTTP, entre outras), que podem ser coletadas e
exibidas em Grafana como no capítulo \ref{cap:obs-rede} \cite{hubbledocs,
prometheusdocs, grafanadocs}. Assim, a cadeia de observabilidade do Manual~1 e do
capítulo \ref{cap:obs-rede} se estende a nível de cluster.

% ----------------------------------------------------------------
\section{Segurança de rede com Cilium}
% ----------------------------------------------------------------

Além das políticas de rede, o Cilium oferece recursos de segurança que se apoiam
no eBPF \cite{ciliumdocs}:

\begin{itemize}
  \item \textbf{Criptografia transparente} entre nós, com WireGuard ou IPsec, ativada
        por parâmetros de instalação, sem alterar as aplicações;
  \item \textbf{Políticas por identidade em vez de endereço}, o que reduz a superfície
        de erro de listas de IP desatualizadas;
  \item \textbf{Visibilidade de descartes} por meio do Hubble, que dá ao administrador
        uma trilha de auditoria de tráfego negado;
  \item \textbf{Integração com o Tetragon} \cite{tetragon}, projeto da mesma
        família, que acrescenta observabilidade e imposição em tempo de execução
        sobre processos e chamadas de sistema.
\end{itemize}

A ativação da criptografia com WireGuard é feita na instalação (comando de cluster,
\textbf{não executado}; os nomes dos parâmetros devem ser conferidos na documentação
da versão):

\begin{lstlisting}[language=bash, caption={Ativação da criptografia transparente com WireGuard (comando de cluster, não executado).}, label=lst:cilium-wg]
cilium install --set encryption.enabled=true --set encryption.type=wireguard
cilium status
\end{lstlisting}

\textbf{Limites e cuidados.} A segurança do Cilium vale tanto quanto a configuração:
uma política \textit{default-allow} (nenhuma política selecionando o \textit{pod})
continua permissiva, e o Hubble só vê o que passa pelos programas do Cilium. O
agente roda com privilégios elevados em cada nó e carrega código no kernel; seu
comprometimento é grave, e o acesso ao cluster deve ser controlado de acordo. Por fim, a
mesma regra de honestidade vale para o restante deste manual: o que foi descrito aqui
representa a documentação oficial \cite{ciliumdocs}, não uma execução. Para validá-lo, o
caminho mais curto é um cluster local de teste (como \textit{kind} ou \textit{minikube}) e o
roteiro da documentação do projeto.
